\documentclass[12pt]{article}
\usepackage[utf8]{inputenc}
\usepackage{amsmath, amssymb, amsthm}
\usepackage{geometry}
\geometry{a4paper, margin=1in}
\usepackage{hyperref}
\usepackage{graphicx}
\usepackage{booktabs}
\usepackage{xcolor}

\newtheorem{proposition}{Proposition}
\newtheorem{remark}{Remark}

\title{\textbf{Part VI: Phenomenological Projections for Extra-Dimensional Geometric Resonance \\[4pt] from K3-Fibered Calabi--Yau Compactifications}}
\author{\textbf{Xavier Callens} \\ \textit{SocrateAI Scientific Agora (independent, non-profit)} \\ \texttt{callensxavier@gmail.com}}
\date{July 2026 \\[0.3cm]
\textbf{DOI:} \href{https://doi.org/10.5281/zenodo.21350629}{10.5281/zenodo.21350629} \\
\textbf{Repository:} \href{https://github.com/xaviercallens/SocrateAI-Scientific-Agora-K3-DarkMatter}{SocrateAI-Scientific-Agora-K3-DarkMatter}}

\begin{document}

\maketitle

\begin{abstract}
We explore the phenomenological consequences of identifying ultralight axion dark matter with the lowest Kaluza--Klein (KK) resonance mode of a $K3$-fibered Calabi--Yau compactification, extending the 5D orbifold dark-photon framework of Lee and Tsai (Phys.\ Rev.\ D, 2026) to an exact 6D string geometry.
We construct a parameterized effective field theory (EFT) in which the KK mass depends on a local baryonic asymmetry parameter $\Delta$ through an exponential warping relation $m_{\mathrm{eff}}(\Delta) = m_0 \exp(k\Delta)$, where $m_0 = 1.83 \times 10^{-21}$~eV and $k$ are \emph{phenomenological} parameters not yet derived from moduli stabilization.
Three sensitivity analyses are presented:
(i)~reprocessing 500 SDSS~DR17 spectroscopic galaxies through this ansatz to project a macroscopic KK resonance map;
(ii)~an achievability demonstration that the Chameleon-type geometric pinching near M87* can shift the effective axion mass above the superradiance exclusion window ($10^{-21}$--$10^{-20}$~eV);
(iii)~a Monte Carlo sensitivity study showing that the ``Cosmic See-Saw'' --- the predicted mass differential between the dense early universe ($z \sim 9$) and the frozen local universe --- would be statistically detectable at $> 5\sigma$ if the underlying density--asymmetry mapping holds.
We emphasize throughout that these results are \emph{illustrative projections} within a parameterized EFT framework, not direct empirical validations.
The free parameters, mock-data assumptions, and open problems requiring top-down string-theoretic input are explicitly enumerated.
All code is publicly available and kernel-verified where applicable.
\end{abstract}

\tableofcontents

%% ====================================================================
\section{Introduction}
\label{sec:intro}

The hypothesis that dark matter (DM) arises as a macroscopic manifestation of particles resonating through a compactified extra dimension has gained renewed theoretical interest with the proposal by Lee and Tsai~\cite{tsai2026} of a 5D $S^1/(Z_2 \times Z_2')$ orbifold model in which a dark photon acquires mass through Kaluza--Klein (KK) quantization.
Their framework demonstrates that the compactification radius $R$ directly dictates the KK mass spectrum,
\begin{equation}
\label{eq:kk_mass_tsai}
m_E^{(n)} = \sqrt{m_B^2 + \frac{n^2}{R^2}}\,,
\end{equation}
where $m_B$ is the brane-localized mass and $n$ indexes the tower of resonances.
This establishes a concrete mechanism by which geometry \emph{creates} mass.

In the present work, we investigate the observational projections that arise when one replaces the Sheffield group's simplified 5D toy geometry with a specific 6D string-theoretic compactification: a $K3 \times T^2$ fibration over a Calabi--Yau threefold, as developed in Parts~I--V of this series~\cite{callens2026partI, callens2026partII}.
The $S_{1,2}$ Picard--Fuchs operator, identified in Part~I as the unique surviving K3 surface candidate under the exact-rational algebraic sieve (with $S_{2,1}$ reclassified as elliptic; see~\cite{callens2026partI}, \S2 and Appendix~A), provides a specific topological rigidity that maps onto the KK resonance structure of Eq.~(\ref{eq:kk_mass_tsai}).

\begin{remark}[Scope and epistemic status]
This paper presents \textbf{phenomenological projections}, not empirical discoveries.
The mapping between the Sheffield 5D framework and our 6D geometry is a physically motivated \emph{ansatz}.
The coupling parameters governing the density-dependent mass variation are free and not derived from moduli stabilization.
We adopt this framing throughout to maintain scientific integrity while exploring the observational consequences of extra-dimensional resonance physics.
\end{remark}

%% ====================================================================
\section{Theoretical Framework}
\label{sec:theory}

\subsection{From 5D Orbifold to 6D K3 Fibration}

The Sheffield model~\cite{tsai2026} compactifies a single extra dimension on $S^1/(Z_2 \times Z_2')$ with radius $R$.
Our framework promotes this to a Type~IIB orientifold on $K3 \times T^2 / \mathbb{Z}_2$, where the $K3$ fiber provides the topological rigidity that fixes the lowest-lying resonance mode, and the $T^2$ torus provides a rolling quintessence modulus (Part~II,~\cite{callens2026partII}).

The $S_{1,2}$ K3 surface is characterized by a third-order Picard--Fuchs recurrence with topological stiffness invariant $V''(0) = 1014$, kernel-verified in Lean~4~\cite{callens2026partI}.
The instanton-generated axion mass is given by the Svrcek--Witten formula~\cite{svrcek2006}:
\begin{equation}
\label{eq:svrcek_witten}
m_a = \frac{M_{\mathrm{Pl}}}{\sqrt{\mathcal{V}}} \sqrt{\sum_d d^2 \, q_d \, e^{-2\pi d \tau}}\,,
\end{equation}
where $\tau$ is the K\"ahler modulus and $\mathcal{V}$ the internal volume.

\begin{remark}[Mass calibration caveat]
\label{rem:mass_caveat}
The axion mass $m_0 \equiv m_a(S_{1,2}) = 3.18 \times 10^{-21}$~eV (and $m_a(S_{2,1}) = 1.83 \times 10^{-21}$~eV) are computed using specific K\"ahler modulus values ($\tau_{12} = 33.6255$, $\tau_{21} = 33.8014$) whose provenance is undocumented --- they reproduce the target masses to 3 significant figures via Eq.~(\ref{eq:svrcek_witten}), raising the possibility of circular calibration.
This is an open provenance question; see Part~I, \S1.2 and \texttt{CAVEATS.md}, \S1.
The masses should be understood as \emph{achievability demonstrations} within a two-parameter $(\tau, \mathcal{V})$ family, not unique predictions.
\end{remark}

\subsection{Density-Dependent Mass Variation: The Chameleon--KK Ansatz}
\label{sec:chameleon_ansatz}

Motivated by the Khoury--Weltman chameleon mechanism~\cite{khoury2004a, khoury2004b} and the density-dependent ``shape-shifting'' proposed by Tsai et al.~\cite{tsai2026}, we parameterize the effective KK resonance mass as:
\begin{equation}
\label{eq:meff}
m_{\mathrm{eff}}(\Delta) = m_0 \, \exp\!\left(k \, \Delta\right)\,,
\end{equation}
where $\Delta \geq 0$ is a dimensionless baryonic asymmetry parameter encoding the local departure from the cosmological mean density, and $k > 0$ is a coupling constant.
An exponential mass scaling is physically well-motivated in string compactifications; it corresponds naturally to a linear backreaction shift in the K\"ahler modulus ($\tau \to \tau - c\Delta$) within the instanton action of Eq. 2.
Physically, increasing $\Delta$ corresponds to the baryonic environment ``pinching'' the K3 extra dimension, shrinking the effective compactification radius $R$ and thereby increasing the KK resonance frequency.

\begin{remark}[Free parameters]
\label{rem:free_params}
Both $k$ and the mapping $\rho_b \mapsto \Delta$ are \textbf{phenomenological}.
In particular:
\begin{itemize}
\item The coupling $k = 0.048$ used in our numerical projections is not derived from K3 moduli stabilization, Type~IIB flux compactification, or any top-down string construction.
\item The standard Khoury--Weltman chameleon scaling $m_{\mathrm{eff}} \propto \rho^{\gamma}$ with $\gamma = 1/4$ maps to an unphysical KW potential index $n = -3$ (see Part~I, \S3 and \texttt{CAVEATS.md}, \S3).
\item A rigorous derivation of $k$ would require the explicit D-brane configuration, moduli stabilization mechanism, and the backreaction of baryonic stress-energy on the K3 metric --- all of which remain open problems (see \texttt{OPEN\_PROBLEMS.md}).
\end{itemize}
This ansatz should therefore be understood as an \emph{effective parameterization} whose top-down embedding is an explicit target for future work.
\end{remark}

%% ====================================================================
\section{Sensitivity Analysis I: Macroscopic KK Resonance Map (WS9)}
\label{sec:ws9}

\subsection{Data and Methodology}

To project how the density-dependent mass ansatz of Eq.~(\ref{eq:meff}) would manifest across a representative galaxy sample, we query the SDSS~DR17 spectroscopic database~\cite{sdss_dr17} for 500 luminous galaxies satisfying:
\begin{itemize}
\item Spectroscopic redshift $0.01 \leq z \leq 0.15$ (local universe);
\item Spectral classification \texttt{GALAXY};
\item Petrosian radius $r_{\mathrm{petro}} > 0$.
\end{itemize}
For each galaxy, we compute a baryonic density proxy:
\begin{equation}
\label{eq:density_proxy}
\tilde{\rho} \equiv \frac{F_r}{r_{\mathrm{petro}}^2}\,, \qquad F_r = 10^{(22.5 - m_r)/2.5}\,,
\end{equation}
where $m_r$ is the extinction-corrected $r$-band magnitude.
The asymmetry parameter is then obtained by linear normalization:
\begin{equation}
\label{eq:delta_normalization}
\Delta_i = \frac{\tilde{\rho}_i}{\max_j \tilde{\rho}_j} \times \Delta_{\max}\,,
\end{equation}
where $\Delta_{\max} = 60$ is chosen to keep $m_{\mathrm{eff}}$ within the astrophysically relevant range.

\begin{remark}[Normalization is arbitrary]
\label{rem:ws9_normalization}
The linear normalization in Eq.~(\ref{eq:delta_normalization}) and the choice $\Delta_{\max} = 60$ are phenomenological prescriptions, not physical derivations.
A different $\Delta_{\max}$ would rescale the entire projected mass map.
This projection demonstrates the \emph{structure} of the density--mass mapping, not its absolute calibration.
We note that $\tilde{\rho}$ represents a 2D projected surface brightness proxy; a rigorous 3D mass-density mapping requires converting optical luminosities to stellar masses ($M_*$) and applying a halo-occupation distribution (HOD) model for 3D deprojection.
\end{remark}

\subsection{Results}

Applying Eq.~(\ref{eq:meff}) with $m_0 = 1.83 \times 10^{-21}$~eV and $k = 0.048$ to the SDSS~DR17 sample yields:
\begin{center}
\begin{tabular}{ll}
\toprule
\textbf{Quantity} & \textbf{Value} \\
\midrule
Mean projected KK mass (local universe) & $1.95 \times 10^{-21}$~eV \\
Peak projected KK mass (highest density) & $3.26 \times 10^{-20}$~eV \\
Galaxies processed & 500 (live SDSS DR17 query) \\
\bottomrule
\end{tabular}
\end{center}
The mean projected mass is close to the base mass $m_0$, reflecting the fact that most galaxies in the local universe reside in relatively low-density environments.
The peak mass, occurring at the densest galaxy core in the sample, is shifted by approximately one order of magnitude --- consistent with the qualitative expectation that baryonic overdensities would drive the KK resonance to higher frequencies.

The full projected map (500 galaxies with coordinates, redshifts, density proxies, and KK masses) is archived at \texttt{data/ws9\_sdss\_kk\_map.csv}.

%% ====================================================================
\section{Sensitivity Analysis II: Superradiance Achievability at M87* (WS10)}
\label{sec:ws10}

\subsection{The Superradiance Constraint}

Ultralight bosons with Compton wavelengths comparable to a black hole's gravitational radius trigger superradiant instabilities that spin down the black hole on astrophysically short timescales~\cite{detweiler1980, dolan2007, arvanitaki2010}.
For M87* ($M = 6.5 \times 10^9 \, M_\odot$; EHT~2019~\cite{eht2019}; $a_* = 0.90$ illustrative, not an EHT measurement), the exclusion window for scalar bosons is approximately $10^{-21}$--$10^{-20}$~eV.

This is precisely the mass range occupied by the $S_{1,2}$ and $S_{2,1}$ axion candidates.
Part~I established via the exact Dolan (2007) continued-fraction method (validated to $<0.4\%$ against published Table~I) that:
\begin{itemize}
\item $S_{2,1}$ ($\alpha_{\mathrm{bare}} = 0.089$): $\tau_{\mathrm{inst}} \approx 380$~Myr $> \tau_{\mathrm{Salpeter}} \approx 50$~Myr --- \textbf{survives unscreened}~\cite{callens2026partI}.
\item $S_{1,2}$ ($\alpha_{\mathrm{bare}} = 0.155$): $\tau_{\mathrm{inst}} \approx 4.6$~Myr $< \tau_{\mathrm{Salpeter}}$ --- \textbf{requires screening}.
\end{itemize}

\subsection{Chameleon Achievability Demonstration}

We ask: \emph{does there exist a physically plausible value of $\Delta$ near M87* such that the Chameleon--KK shift of Eq.~(\ref{eq:meff}) moves $S_{1,2}$ above the exclusion window?}

Setting $\Delta_{\mathrm{M87*}} = 150$ (intended to represent the extreme baryonic overdensity of the accretion environment), Eq.~(\ref{eq:meff}) yields:
\begin{equation}
m_{\mathrm{eff}}(\Delta = 150) = 1.83 \times 10^{-21} \times e^{0.048 \times 150} \approx 2.45 \times 10^{-18}~\text{eV}\,,
\end{equation}
which lies approximately two orders of magnitude above the superradiance exclusion window.

\begin{remark}[Circularity disclosure]
\label{rem:ws10_circularity}
The value $\Delta = 150$ is \textbf{not computed} from the gas density profile, accretion rate, or GRMHD modeling of M87*.
It is chosen to demonstrate that the exponential form of Eq.~(\ref{eq:meff}) can, in principle, produce sufficient mass uplift.
This is an \emph{achievability check} --- it establishes that the Chameleon--KK mechanism is not trivially excluded, but it does not constitute a prediction that M87*'s environment actually induces $\Delta = 150$.
A genuine validation would require computing $\Delta$ from first principles using the baryonic stress-energy tensor of a realistic accretion flow and its backreaction on the K3 metric, which is beyond the scope of this work.
\end{remark}

%% ====================================================================
\section{Sensitivity Analysis III: Cosmic See-Saw Detectability (WS11)}
\label{sec:ws11}

\subsection{The See-Saw Hypothesis}

The ``Cosmic See-Saw'' posits that the K3 geometric resonance was strongly active in the dense early universe ($z \gtrsim 9$), producing heavier effective axion masses that accelerated structure formation, but has ``frozen out'' in the dilute local universe ($z \approx 0$), where the $S_{1,2}$ geometry has settled into its perturbative ground state ($S_{1,2} \leq 1.177$ --- a bound reported by the companion \texttt{DarkMatterK3} GPU pipeline from 327,918 SDSS DR17 galaxies; the converged run behind it is external to this repository's committed artifacts, and the bound carries the epistemic status of a reported pipeline output pending an archived reproduction script; see Paper~V, \S``Empirical Status, Provenance, and Physical Interpretation'').
This qualitative picture is consistent with recent JWST UNCOVER observations of unexpectedly massive galaxies at $z \sim 9$--$11$ and the mature structure of the local cosmic web.

\subsection{Monte Carlo Sensitivity Study}

To quantify whether such a mass differential would be statistically detectable, we construct synthetic $\Delta$ distributions designed to represent the two epochs:
\begin{align}
\Delta_{\mathrm{local}} &\sim \mathcal{N}(\mu = 5, \sigma = 2), \quad N = 1000 \quad \text{(local freeze-out)}\,, \label{eq:delta_local} \\
\Delta_{\mathrm{early}} &\sim \mathcal{N}(\mu = 45, \sigma = 10), \quad N = 1000 \quad \text{(early-universe active resonance)}\,. \label{eq:delta_early}
\end{align}
These choices reflect the qualitative astrophysical expectation that early-universe structures are $\sim 10\times$ denser than local ones, but they are \textbf{not fits to observational data}.

\begin{remark}[Mock data disclosure]
\label{rem:ws11_mock}
The distributions in Eqs.~(\ref{eq:delta_local})--(\ref{eq:delta_early}) are \textbf{entirely synthetic}.
They are constructed to model the qualitative density contrast between the JWST UNCOVER epoch and the local SDSS universe, not derived from photometric or spectroscopic measurements of individual high-$z$ galaxies.
Because the two distributions have well-separated means, a Welch $t$-test is \emph{guaranteed} to return an extremely low $p$-value.
This analysis therefore quantifies the \emph{projected statistical power} of a See-Saw detection given a hypothesized density contrast --- it does not constitute an empirical validation of the See-Saw hypothesis.
\end{remark}

\subsection{Results}

Applying Eq.~(\ref{eq:meff}) to both synthetic samples and performing a two-sided Welch's $t$-test on $\log_{10}(m_{\mathrm{eff}})$:
\begin{center}
\begin{tabular}{ll}
\toprule
\textbf{Statistic} & \textbf{Value} \\
\midrule
Mean $m_{\mathrm{eff}}$ (local, $z \approx 0$) & $2.34 \times 10^{-21}$~eV \\
Mean $m_{\mathrm{eff}}$ (early, $z \approx 9$) & $1.84 \times 10^{-20}$~eV \\
Welch $t$-statistic & 126.53 \\
$p$-value & $< 10^{-300}$ (numerically zero) \\
Projected confidence & $> 10\sigma$ \\
\bottomrule
\end{tabular}
\end{center}

\noindent \textbf{Interpretation:}
If the density--asymmetry mapping of Eq.~(\ref{eq:meff}) is correct and the true $\Delta$ distributions at $z \sim 0$ and $z \sim 9$ are qualitatively as modeled, the Cosmic See-Saw signature would be overwhelmingly detectable.
This establishes the See-Saw as a \emph{falsifiable prediction} of the K3 resonance framework: future surveys (Euclid, Rubin LSST) mapping baryonic density across redshift could test whether the required density contrast exists and whether it correlates with dark matter phenomenology as predicted.

%% ====================================================================
\section{Relationship to the Sheffield Framework}
\label{sec:sheffield}

Table~\ref{tab:sheffield_comparison} summarizes the conceptual mapping between the Tsai et al.\ (2026) 5D framework and the K3 6D extension.
The division of labor is deliberate: the $S^1/(Z_2 \times Z_2')$ orbifold supplies the \emph{microscopic} mass-generation mechanism (KK quantization against the radius $R$), while the TDA pipeline supplies the \emph{macroscopic} observable ($\Delta$, with the local baryonic field pinching the effective $R$).
No simulation of the orbifold itself is required on the observatory side: the 5D physics enters only through the parameterized response $m_{\mathrm{eff}}(\Delta)$ of Eq.~(\ref{eq:meff}).
A fuller discussion of how the reported local-universe bound $S_{1,2} \le 1.177$, the sector-level stability ($S_{1,2} < 1.3$ everywhere vs.\ the model-internal transition threshold $1.8$), and the Poisson-mock null tests relate to the Sheffield freeze-out requirement --- including the provenance limits of each number --- is given in Paper~V, \S``Empirical Status, Provenance, and Physical Interpretation.''

\begin{table}[h]
\centering
\caption{Conceptual bridge: Sheffield 5D $\to$ K3 6D.}
\label{tab:sheffield_comparison}
\begin{tabular}{p{3cm}p{5cm}p{5cm}}
\toprule
\textbf{Mechanism} & \textbf{Sheffield 5D} & \textbf{K3 $\times$ $T^2$ 6D (this work)} \\
\midrule
Mass generation & KK quantization on $S^1/(Z_2 \times Z_2')$ & Instanton action on $S_{1,2}$ K3 fiber \\
Density dependence & Warped extra dimension & Chameleon--KK ansatz, Eq.~(\ref{eq:meff}) \\
Early/late dichotomy & Active resonance $\to$ freeze-out & Cosmic See-Saw (\S\ref{sec:ws11}) \\
Observational tool & Not specified & Projected SDSS~DR17 density map (\S\ref{sec:ws9}; Euclid extension planned, no Euclid data processed to date) \\
Free parameters & $R$, brane mass $m_B$ & $k$, $\Delta_{\max}$, $(\tau, \mathcal{V})$ \\
\bottomrule
\end{tabular}
\end{table}

%% ====================================================================
\section{Open Problems and Required Inputs}
\label{sec:open_problems}

The following problems must be resolved before the phenomenological projections of this paper can be promoted to quantitative predictions.
We publish them here to establish priority and to invite collaboration from the string phenomenology and observational cosmology communities.

\begin{enumerate}
\item \textbf{Top-down derivation of $k$.} The coupling constant in Eq.~(\ref{eq:meff}) must be derived from the backreaction of baryonic stress-energy on the K3 metric within a stabilized Type~IIB compactification.
This requires the explicit D-brane configuration, which is not established (see \texttt{OPEN\_PROBLEMS.md}).

\item \textbf{$\Delta$ from GRMHD.} The baryonic asymmetry parameter near M87* must be computed from realistic accretion-flow simulations (e.g., the EHT GRMHD library), not assumed.

\item \textbf{$S_{2,1}$ reclassification.} While $S_{2,1}$ fails the strict K3 order-3 classification, its order-2 (elliptic) nature natively embeds into standard F-theory compactifications on elliptically fibered Calabi-Yau fourfolds. Future work will assess whether $S_{2,1}$ survives as a valid axion generator in this alternative F-theory context.

\item \textbf{Chameleon $\gamma$ tension.} The fitted $\gamma = 0.25$ maps to unphysical KW index $n = -3$.
A physically admissible screening mechanism (symmetron, native-modulus, or $\gamma = 1/2$ with boosted density ratio) must be identified.

\item \textbf{Boltzmann-level cosmology.} The $H_0$ and $S_8$ predictions require a genuine CLASS C-source fork, not background-only integration (see Part~II, \S5 and \texttt{CAVEATS.md}, \S5).
The self-consistent background recomputation yields $H_0 \approx 75.8$~km/s/Mpc, not the originally reported 71.92.

\item \textbf{JWST UNCOVER fit.} The early-universe ($z \sim 9$) $\Delta$ distribution (Eq.~\ref{eq:delta_early}) must be replaced with a fit to actual JWST UNCOVER photometric/spectroscopic data.
\end{enumerate}

%% ====================================================================
\section{Conclusions}
\label{sec:conclusions}

We have presented three phenomenological sensitivity analyses exploring the observational consequences of identifying ultralight axion dark matter with the Kaluza--Klein resonance of a K3-fibered Calabi--Yau compactification, extending the recent Sheffield 5D framework to a concrete 6D string geometry.

The results demonstrate that:
\begin{itemize}
\item The density-dependent mass ansatz $m_{\mathrm{eff}} = m_0 \exp(k\Delta)$ produces a physically structured macroscopic mass map when applied to real SDSS~DR17 galaxy data (\S\ref{sec:ws9}).
\item The Chameleon--KK mechanism can, in principle, shift the axion mass sufficiently to evade M87* superradiance bounds, though the required environmental parameters are not yet derived from first principles (\S\ref{sec:ws10}).
\item The Cosmic See-Saw signature would be overwhelmingly detectable ($> 10\sigma$) if the hypothesized density contrast between the early and local universe maps onto the assumed $\Delta$ distributions (\S\ref{sec:ws11}).
\end{itemize}

We emphasize that these are \emph{projections within a parameterized EFT}, not empirical discoveries.
The free parameters $(k, \Delta_{\max}, \gamma)$ and the mock distributions (Eqs.~\ref{eq:delta_local}--\ref{eq:delta_early}) are explicitly disclosed.
Promoting these projections to predictions requires the top-down string-theoretic inputs enumerated in \S\ref{sec:open_problems}.

All code, data, and Lean~4 formal proofs are publicly available at the repository listed on the title page.

%% ====================================================================
\section*{Acknowledgments}

The author thanks Taegyu Lee and Yu-Dai Tsai for the theoretical inspiration provided by their extra-dimensional resonance framework, and the SDSS~DR17 collaboration for public data access.
This work was conducted within the SocrateAI Scientific Agora framework.
The author is an independent researcher and this work is not affiliated with any university or national laboratory.

%% ====================================================================
\begin{thebibliography}{99}

\bibitem{tsai2026}
T.~Lee and Y.-D.~Tsai, ``Naturally resonant dark matter from extra dimensions,''
\textit{Phys.\ Rev.\ D} \textbf{114}, L011701 (2026).

\bibitem{callens2026partI}
X.~Callens, ``Part I: String-Inspired Effective Field Theories from K3 Surfaces: Resolving Fuzzy Dark Matter Tensions via Exact Algebraic Sieving,''
Zenodo, \href{https://doi.org/10.5281/zenodo.21350629}{10.5281/zenodo.21350629} (2026).

\bibitem{callens2026partII}
X.~Callens, ``Part II: Project Vafa-Continuity: A String-Inspired $K3 \times T^2$ Quintessence Model and the Swampland Obstruction to Stable Dark Energy,''
Zenodo, \href{https://doi.org/10.5281/zenodo.21350629}{10.5281/zenodo.21350629} (2026).

\bibitem{svrcek2006}
P.~Svrcek and E.~Witten, ``Axions in string theory,''
\textit{JHEP} \textbf{0606}, 051 (2006);
\href{https://arxiv.org/abs/hep-th/0605206}{hep-th/0605206}.

\bibitem{khoury2004a}
J.~Khoury and A.~Weltman, ``Chameleon fields: Awaiting surprises for tests of gravity in space,''
\textit{Phys.\ Rev.\ Lett.} \textbf{93}, 171104 (2004).

\bibitem{khoury2004b}
J.~Khoury and A.~Weltman, ``Chameleon cosmology,''
\textit{Phys.\ Rev.\ D} \textbf{69}, 044026 (2004).

\bibitem{sdss_dr17}
SDSS Collaboration, ``Sloan Digital Sky Survey Data Release 17,''
\textit{ApJS} \textbf{259}, 35 (2022).

\bibitem{detweiler1980}
S.~Detweiler, ``Klein-Gordon equation and rotating black holes,''
\textit{Phys.\ Rev.\ D} \textbf{22}, 2323 (1980).

\bibitem{dolan2007}
S.~R.~Dolan, ``Instability of the massive Klein-Gordon field on the Kerr spacetime,''
\textit{Phys.\ Rev.\ D} \textbf{76}, 084001 (2007).

\bibitem{arvanitaki2010}
A.~Arvanitaki, S.~Dimopoulos, S.~Dubovsky, N.~Kaloper, and J.~March-Russell,
``String axiverse,''
\textit{Phys.\ Rev.\ D} \textbf{81}, 123530 (2010).

\bibitem{eht2019}
Event Horizon Telescope Collaboration, ``First M87 Event Horizon Telescope Results. I.,''
\textit{ApJL} \textbf{875}, L1 (2019).

\bibitem{planck2018}
Planck Collaboration, ``Planck 2018 results. VI. Cosmological parameters,''
\textit{A\&A} \textbf{641}, A6 (2020).

\end{thebibliography}

\end{document}
